%% title: The most likely properties of the labelling algorithm
%% category: misc

\section{Hypothesized properties}

As an initial hypothesis to be modified, I propose that the labelling algorithm (LA) underlying human syntax has the following properties:

% what MAKES associativity in a system?

\begin{enumerate}
    \item \textbf{Not adjunct-reading.} The LA may not refer to an adjunct's features. This will be assumed throughout definitions of the other properties. 

    \item \textbf{Nearsighted.} The LA cannot look through pointers between syntactic nodes: it can only see the features of the nodes directly inputted into it, not the features of their daughters. In other words, the LA's input consists of two \texttt{snode.values{}} objects, not two \texttt{snode} objects.

    \item \textbf{Noncommutative.} The LA receives ordered inputs. The order need not matter intrinsically; noncommutativity only mandates that $\tau (x \magma y) = \tau(y \magma x)$ is not always true. 

    \item \textbf{Wholesighted.} The LA can distinguish between syntactic objects which have complete and incomplete semantics: in other words, between maximal projections (specifiers and complements), and less-than-maximal projections (heads and bars). This is a means by which the LA can be noncommutative.\footnote{The alternative means for noncommutativity include distinguishing between left and right daughters (specifiers and heads on the one hand, and complements and bars on the other), or making a tripartite distinction between specifiers, complements, and heads. Distinguishing between minimal and non-minimal projections seems less principled, as the LA would need to know how many semantic arguments every head takes (assuming heads take multiple semantic arguments).}

    \item \textbf{Headed-sequential.} The LA may not contain rules which jointly refer to specifiers and complements: a language-specific LA can be specified with only rules of the form $\tau(\text{head}, \text{compl})$ and $\tau(\text{spec}, \text{bar})$.\footnote{The alternatives are \textbf{Decapitated-sequential}, where the LA first operates over a spec-compl merge, and then over their head, and \textbf{Nonsequential}, where the LA can condition on specifiers, heads, and complements simultaneously.}

    \item \textbf{Tier-blind.} For each pair of feature values in a given LA, there is only one possible mother value: a labelling operation over tier X (e.g. number) cannot make reference to features on tier Y (e.g. person). 

    \item \textbf{Conservative.} For each mother feature, its value is one of the daughter values. This is a crucial property which severely restricts the possibility space for agreement patterns. 

    \item \textbf{Specifier-reading.} The LA may refer to a specifier's features: Spec-Head agreement is possible.

    \item \textbf{Associative.} The mother feature of [X[YZ]] will always be the same as the mother feature of [[XY]Z]: there may be no `rock paper scissors' configurations where value X beats value Y, value Y beats value Z, and value Z beats value X. 

    \item \textbf{Value-seeking.} If at least one readable daughter is valued for a feature, its mother may not be unvalued for that feature. (Therefore, the LA must be surjective.)

    \item \textbf{Consistently acommutative.}

    \item \textbf{Superconsistently acommutative.}

    \item \textbf{Whole-seeking.}

    \item \textbf{Partial-seeking.}

    \item \textbf{Bare phrase structure.} Maximal projections and minimal projections are underlyingly the same object. Therefore, features that the maximal projection inherits are also inherited by the minimal projection.\footnote{Maybe this is eclipsed by mother-sharing.}

    % \item \textbf{Whole-writing.} After the features of a mother node are calculated, an algorithm runs which potentially donates some of them to its whole daughter (specifier or complement). 

    % \item \textbf{Non-overriding.} A mother may only value unvalued features of its whole complement -- it cannot override existing valued features. (This is something like the Activity Condition.)\footnote{I'm really unsure about this one.}

    % \item \textbf{Not partial-writing.} Mother nodes cannot donate features to their non-whole/`partial' daughters (heads and bars), beyond the fact that maximal and minimal projections are the same objects. Partial-writing would give rise to a recursive process through which all specifiers could pass features to complements, or even farther down the clausal spine.

    \item \textbf{Mother-sharing.} When a phase is spelled out, starting at the top of the phase and proceeding downwards, mothers share features with their children. 
        
    % \item \textbf{Phase-sharing.} When a phase is spelled out, it shares its features to all nodes lower down in the tree recursively. 

    \item \textbf{Non-overriding.} During mother-sharing, only unvalued features may be valued. Valued features may not be overridden/re-valued. 


    % Question: is there something about category features, in that we have:
    % - case, where potentially category features of functional heads get attached to wholes and don't trigger anti-overriding
    % - allomorphy and allosemy, where things can see category features (ok maybe other things too tho) of nodes around them?

    % maybe case is a special 'phase' feature tho actually! or 'domain' or 'semantic part of things' or smth
    % which only phase heads normally have



\end{enumerate}

\begin{tryit}{report}{[CP* C [TP T(number=?pl) [vP* [DP D(number=pl)] v [VP V [DP D(number=sg) [NP N +[AP A(number=?sg)]]]]]]]}
Type a tree into the box below in the following format. If the tree passes, that means it can be represented by a consistent labelling algorithm with the properties outlined above, in which case the mandated rules of the labelling algorithm are outputted. 

  [XP ...]          a phrase; write XP* if it is a phase
  spec X compl      any number of specifiers, then the head, then at most one complement
  +[AP ...]         an adjunct
  X(tier=value)     a head with a lexical/inherent feature value
  X(tier=?)         a head unvalued for that tier
  X(tier=?value)    unvalued, and observed to surface with a specific value
  X(tier=?-)        unvalued, and observed to stay unvalued
\end{tryit}

\begin{minted}{python}
import itertools
import re
import sys


# define the basic data structures
class Head:
    def __init__(self, name, lexical, observed):
        self.name = name
        self.lexical = lexical
        self.observed = observed

class Phrase:
    def __init__(self, name, phase, head, complement, specifiers, adjuncts):
        self.name = name
        self.phase = phase
        self.head = head
        self.complement = complement
        self.specifiers = specifiers
        self.adjuncts = adjuncts

    # return all daughters in a dependency tree
    def daughters(self):
        return [d for d in (self.complement, *self.specifiers, *self.adjuncts) if d is not None]

    # recursively yield all subphrases, including self
    def phrases(self):
        yield self
        for d in self.daughters():
            yield from d.phrases()

# parse the inputted tree string into objects
TOKEN = re.compile(r"\[|\]|\+|[^\s\[\]()+]+(?:\([^)]*\))?")

# parse a head object into a head data structure
def parse_head(token):
    m = re.fullmatch(r"([^()]+)(?:\((.*)\))?", token)
    if not m:
        raise ValueError(f"cannot read head {token!r}")
    lexical, observed = {}, {}
    for feature in filter(None, (f.strip() for f in (m.group(2) or "").split(","))):
        if "=" not in feature:
            raise ValueError(f"feature {feature!r} in {token!r} needs the form tier=value")
        tier, value = (s.strip() for s in feature.split("=", 1))
        if not value.startswith("?"):
            lexical[tier] = value
        elif value == "?":
            lexical[tier] = None
        else:
            lexical[tier] = None
            observed[tier] = None if value == "?-" else value[1:]
    return Head(m.group(1), lexical, observed)

# parse the whole string of objects into a tree
def parse(text):
    tokens = TOKEN.findall(text)
    pos = 0

    def peek():
        if pos >= len(tokens):
            raise ValueError("the tree ends too early: a ] is missing")
        return tokens[pos]

    def phrase():
        nonlocal pos
        if peek() != "[":
            raise ValueError(f"expected [ but found {peek()!r}")
        pos += 1
        name = peek()
        if name in "[]+":
            raise ValueError("every phrase needs a name, as in [DP ...]")
        pos += 1
        before, head, after, adjuncts = [], None, [], []
        while peek() != "]":
            if peek() == "+":
                pos += 1
                adjuncts.append(phrase())
            elif peek() == "[":
                (after if head else before).append(phrase())
            else:
                if head:
                    raise ValueError(f"{name} has two heads: {head.name} and {peek()}")
                head = parse_head(peek())
                pos += 1
        pos += 1
        if head is None:
            raise ValueError(f"{name} has no head")
        if len(after) > 1:
            raise ValueError(f"{name} has more than one complement")
        return Phrase(name.rstrip("*"), name.endswith("*"), head, after[0] if after else None,
                      list(reversed(before)), adjuncts)

    tree = phrase()
    if pos != len(tokens):
        raise ValueError(f"unexpected {tokens[pos]!r} after the end of the tree")
    return tree

# data structure for a language's transition function
# this is important! with the above properties, a transition
# function can be exhaustively described by (1) a ranking
# of values on that tier, including ties, and (2) a ranking
# of partial (head/bar) vs whole (complement/specifier).
# When two values are tied, they get evaluated according to
# the partial/whole ranking. 
# The property of "Consistently acommutative" describes a grammar
# in which the partial/whole ranking is consistent across all
# value ties on a tier, and the property of "Superconsistently
# acommutative" describes a grammar in which this ranking is
# consistent across all tiers. If it turns out that partials
# always rank over wholes or vice versa, across languages, 
# then the labelling algorithm for every language will either be
# Whole-seeking or Partial-seeking. 
# In this implementation, I assume it is only necessarily
# consistently acommutative. 
class Grammar:
    def __init__(self, blocks, side):
        self.rank = {v: i for i, block in enumerate(blocks) for v in block}
        self.blocks = blocks
        self.side = side

    def __call__(self, partial, whole):
        if partial is None or whole is None:
            return whole if partial is None else partial
        if self.rank[partial] != self.rank[whole]:
            return max(partial, whole, key=self.rank.get)
        return partial if self.side == "partial" else whole

# checks that every transition rule is conservative, 
# value-seeking, and associative
def lawful(tau, values):
    domain = (None, *values)
    for x in domain:
        for y in domain:
            m = tau(x, y)
            if m not in (x, y):
                return False
            if m is None and (x, y) != (None, None):
                return False
            if any(tau(tau(x, y), z) != tau(x, tau(y, z)) for z in domain):
                return False
    return True

# lists every possible way to rank the values of a tier incl. ties
def ordered_partitions(items):
    if not items:
        yield []
        return
    for k in range(1, len(items) + 1):
        for block in itertools.combinations(items, k):
            rest = [x for x in items if x not in block]
            for tail in ordered_partitions(rest):
                yield [block, *tail]

# then expands tied values to partial>whole and whole>partial
def grammars(values):
    for blocks in ordered_partitions(sorted(values)):
        sides = ("partial", "whole") if any(len(b) > 1 for b in blocks) else ("partial",)
        for side in sides:
            tau = Grammar(blocks, side)
            if lawful(tau, values):
                yield tau

# runs one tier of one grammar over the tree
def derive(tree, tier, tau):
    head, bars = {}, {}

    # bottom-up labelling pass:
        # first recursively labels everything lower
        # when labelling, first labels nodes with their lexical features
        # then merges the complement's label
        # then merges the result with each specifier in turn,
        # going from innermost to outermost specifier
    def label(ph):
        for d in ph.daughters():
            label(d)
        head[ph] = ph.head.lexical.get(tier)
        chain = [tau(head[ph], bars[ph.complement][-1] if ph.complement else None)]
        head[ph] = chain[-1]
        for spec in ph.specifiers:
            chain.append(tau(chain[-1], bars[spec][-1]))
            head[ph] = chain[-1]
        bars[ph] = chain
        if ph.phase:
            share(ph)

    # how each daughter receives value from mother in mother-sharing
    def inherit(value, ph):
        if bars[ph][-1] is None and value is not None:
            bars[ph][-1] = value
            if head[ph] is None:
                head[ph] = value
        if not ph.phase:
            share(ph)

    # mother-sharing/top-down pass
    def share(ph):
        chain = bars[ph]
        for i in range(len(chain) - 1, 0, -1):
            if chain[i - 1] is None:
                chain[i - 1] = chain[i]
            inherit(chain[i], ph.specifiers[i - 1])
        if head[ph] is None:
            head[ph] = chain[0]
        if ph.complement is not None:
            inherit(chain[0], ph.complement)
        for adjunct in ph.adjuncts:
            inherit(chain[-1], adjunct)

    label(tree)
    if not tree.phase:
        share(tree)
    return {ph.head: head[ph] for ph in tree.phrases()}

# collect all consistent grammars, find mandated rules, 
def check_tier(tree, tier):
    heads = [ph.head for ph in tree.phrases()]
    values = {h.lexical.get(tier) for h in heads} | {h.observed[tier] for h in heads if tier in h.observed}
    values.discard(None)
    targets = [h for h in heads if tier in h.observed]
    passing, possible, total = [], {id(h): set() for h in targets}, 0
    for tau in grammars(values):
        total += 1
        surface = derive(tree, tier, tau)
        for h in targets:
            possible[id(h)].add(surface[h])
        if all(surface[h] == h.observed[tier] for h in targets):
            passing.append(tau)
    pairs = [(p, w) for p in sorted(values) for w in sorted(values) if p != w]
    mandated = [(p, w, passing[0](p, w)) for p in sorted(values) for w in sorted(values)
                if p != w and passing and len({tau(p, w) for tau in passing}) == 1]
    blocked = [(h, sorted(possible[id(h)], key=str)) for h in targets
               if h.observed[tier] not in possible[id(h)]]
    return passing, total, mandated, len(pairs), blocked


def show(value):
    return "unvalued" if value is None else value


def report(text):
    tree = parse(text)
    tiers = sorted({t for ph in tree.phrases() for t in ph.head.lexical})
    if not tiers:
        return "No features on any head, so there is nothing to check."
    lines, ok = [], True
    for tier in tiers:
        passing, total, mandated, pairs, blocked = check_tier(tree, tier)
        if passing:
            lines.append(f"{tier}: passes ({len(passing)} of {total} lawful transition functions derive it)")
            for p, w, m in mandated:
                lines.append(f"  mandated: tau(partial={p}, whole={w}) = {m}")
            if not mandated and pairs:
                lines.append("  no transition rules are mandated")
        else:
            ok = False
            lines.append(f"{tier}: FAILS (none of {total} lawful transition functions derive it)")
            for h, values in blocked:
                lines.append(f"  {h.name} never surfaces as {show(h.observed[tier])}; it can only surface as "
                             + ", ".join(show(v) for v in values))
            if not blocked:
                lines.append("  each observation is derivable on its own, but not all of them together")
    lines.insert(0, "PASSES" if ok else "FAILS")
    return "\n".join(lines)


def main():
    if len(sys.argv) > 1:
        print(report(" ".join(sys.argv[1:])))
        return
    while True:
        try:
            line = input("tree> ")
        except EOFError:
            break
        if not line.strip():
            break
        try:
            print(report(line))
        except ValueError as error:
            print("error:", error)
        print()


if __name__ == "__main__":
    main()
\end{minted}
